\section{Introduction}
\label{sec:introduction}

An inference optimization can preserve the next token while changing the scores or key--value (KV) state that produced it. Consumers that inspect complete scores, replay decisions, or resume saved state depend on this richer numerical interface. Preserving an established reference lets execution change without changing that interface. Deterministic kernels and stateful inference systems already demonstrate such freedom under numerical constraints~\cite{thinkingmachines,llm42,vosti,detent}. We study how to retain useful parallel work inside a quantized decoder once its reference is fixed.

Prefill and speculative verification make several queries known in advance. Their projections can share weight access and activation preparation when quantization, reduction, and rounding remain compatible with the reference. Attention has a separate requirement: a query may need a particular physical KV extent and layout even when its causal mask preserves logical visibility. A legal projection group can therefore cross a boundary at which attention needs a different physical view. Splitting the complete decoder at every such boundary preserves the views but also fragments compatible projection and feed-forward work.

Consider four known queries with reference attention extents $[256,256,256,512]$. An admitted four-column projection processes them together, while attention needs a three-query view and a one-query view. Whole-decoder calls of three and one query impose that finer partition on all surrounding operators. The decoder could instead retain the four-query group and refine only the attention operation that needs the split.

We implement this \emph{local lowering}: form the required Q/K/V and mask views inside attention, execute its admitted tiles, and join the outputs in logical query order. Surrounding operators keep their admitted groups. A stateful consumer also needs the correct output row after rejection and the correct state for its next call. The frontend owns requested vocabulary rows before native storage is reused, commits the accepted prefix and pending token, and checks the successor's complete attention-view identity before graph reuse. Local execution choices thus preserve both the outputs a consumer reads and the state from which decoding continues.

We evaluate the construction through three distinct questions. Direct LOCAL/SIMPLE transactions test whether refinement preserves useful work across a physical crossing and returns reference-equivalent outputs and state. Independent complete-request studies quantify the effects of few-column projection arithmetic (C4) and packing multiple admitted groups in prefill (P4) under their own controls. Natural workloads then reveal how often the relevant geometry occurs and how sampling, drafting, and other request costs affect the benefit. The selected crossing reduces verification-through-commit time by 21.56\%; the independent RTX few-column comparison reduces complete-request time by 15.80\% relative to grouped scalar execution. Their designs and statistical roles are specified in Section~\ref{sec:design}.

The paper makes three contributions:
\begin{enumerate}
\item We identify the conflict between projection arithmetic groups and physical attention partitions through distinct numerical counterexamples (Section~\ref{sec:reference}).
\item We implement a decoder lowering that preserves compatible groups across finer attention boundaries and carries ordered outputs through ownership, commit, and guarded successor reuse (Sections~\ref{sec:method}--\ref{sec:implementation}). It builds on existing numerical kernels and checks complete vectors and active state against their references.
\item We measure the construction's direct transaction effect and separately evaluate few-column arithmetic and grouped prefill over complete requests. Natural-workload and cost studies explain the deployment conditions under which this retained work benefits service (Sections~\ref{sec:design}--\ref{sec:discussion}).
\end{enumerate}
